% 한글 PDF 빌드 header (pandoc -H). tectonic(XeTeX) + xeCJK + tcolorbox.
% 💡 등 컬러 이모지는 라틴/CJK 폰트에 글리프가 없어 newunicodechar 로 단색 대체.
\usepackage{newunicodechar}
\usepackage{fontawesome5}
\usepackage{amssymb}
% --- 이모지/기호 → 렌더 가능한 단색 글리프 매핑 ---
\newunicodechar{💡}{\textcolor[rgb]{0.85,0.5,0}{\faLightbulb}}
\newunicodechar{🦟}{\textbf{*}}
\newunicodechar{✅}{\textcolor[rgb]{0.1,0.6,0.1}{\checkmark}}
\newunicodechar{❌}{\textcolor[rgb]{0.8,0.1,0.1}{\texttimes}}
\newunicodechar{⚠}{\textbf{!}}
\newunicodechar{⚡}{\faBolt}
\newunicodechar{★}{\textbf{*}}
\newunicodechar{☆}{\textbf{*}}
\newunicodechar{✓}{\checkmark}
\newunicodechar{✔}{\checkmark}
\newunicodechar{✗}{\texttimes}
\newunicodechar{✘}{\texttimes}
\newunicodechar{→}{$\rightarrow$}
\newunicodechar{←}{$\leftarrow$}
\newunicodechar{⇒}{$\Rightarrow$}
\newunicodechar{≈}{$\approx$}
\newunicodechar{≤}{$\leq$}
\newunicodechar{≥}{$\geq$}
\newunicodechar{×}{$\times$}
\newunicodechar{±}{$\pm$}
\newunicodechar{•}{\textbullet}
\newunicodechar{◆}{\textbf{*}}
\newunicodechar{▶}{\textbf{>}}
\newunicodechar{−}{\ensuremath{-}}
\newunicodechar{∈}{\ensuremath{\in}}
\newunicodechar{⊙}{\ensuremath{\odot}}
\newunicodechar{≠}{\ensuremath{\neq}}
\newunicodechar{≪}{\ensuremath{\ll}}
\newunicodechar{≫}{\ensuremath{\gg}}
\newunicodechar{↔}{\ensuremath{\leftrightarrow}}
\newunicodechar{↑}{\ensuremath{\uparrow}}
\newunicodechar{↓}{\ensuremath{\downarrow}}
\newunicodechar{※}{\textasteriskcentered}
\newunicodechar{○}{\ensuremath{\circ}}
\newunicodechar{◎}{\ensuremath{\circledcirc}}
\newunicodechar{△}{\ensuremath{\triangle}}
\newunicodechar{①}{(1)}
\newunicodechar{②}{(2)}
\newunicodechar{③}{(3)}
\newunicodechar{④}{(4)}
\newunicodechar{⑤}{(5)}
\newunicodechar{☐}{$\square$}
\newunicodechar{☑}{\textbf{[v]}}
\newunicodechar{️}{}
% 누락 글리프는 경고만 내고 진행 (빌드 중단 방지)
\tracinglostchars=1
% --- blockquote 를 콜아웃 박스로 (💡 블록 강조) ---
\usepackage[most]{tcolorbox}
\tcbuselibrary{breakable}
\renewenvironment{quote}{%
  \begin{tcolorbox}[breakable,colback=orange!6,colframe=orange!55,boxrule=0.5pt,arc=2pt,left=5pt,right=5pt,top=3pt,bottom=3pt]%
}{\end{tcolorbox}}
% 표가 페이지를 넘어가도 깨지지 않게
\usepackage{longtable}
